\subsection{Exercises Series 5}
\vspace*{0.2cm}
\noindent\Large{\underline{\textbf{Exercise 1}}}
\normalsize{}
\begin{enumerate}
    \item Write an equivalent series with the index of summation beginning at $n = 1$ :
    $$  \sum_{n \in \mathbb{N}} \dfrac{x^{2n+1}}{(2n+1)!}, \qquad \sum_{n \in \mathbb{N}^*} (-1)^n (n+1)\, x^n, \qquad 
        \sum_{n \in \mathbb{N}} \dfrac{(-1)^n}{2n+1} x^{2n+1}.  $$
    \item Determine the radius of convergence of the following power series:
    $$  \sum_{n \in \mathbb{N}} (\sin n)\, x^n, \qquad \sum_{n \in \mathbb{N}} \dfrac{x^n}{(n+1)^{a+1} 3^n}, \qquad 
        \sum_{n \in \mathbb{N}} \arccos\!\left(1 - \dfrac{1}{n^2}\right)x^n, \qquad 
        \sum_{n \in \mathbb{N}} \dfrac{n^n}{n!}\, x^n. $$
    \item Find the interval of convergence of the following power series
    $$  \sum_{n \in \mathbb{N}} \dfrac{(-1)^n}{2^n} (x+1)^n, \qquad 
        \sum_{n \in \mathbb{N}^*} \dfrac{x^n}{n}, \qquad \sum_{n \in \mathbb{N}} \dfrac{(-1)^n\, n!}{3^n} (x-5)^n. $$
\end{enumerate}

\noindent\Large{\underline{\textbf{Exercise 2}}}
\normalsize{}\\
\noindent Determine the radius of convergence as well as the sum of the following power series:
    $$ \sum_{n \in \mathbb{N}} \dfrac{x^n}{(n+1)(n+3)}, \quad 
       \sum_{n \in \mathbb{N}} \left(n^2 - n - 3\right) 3^{n-1} x^n, \quad 
       \sum_{n \in \mathbb{N}} \cosh(na)\, x^n, \quad 
       \sum_{n \in \mathbb{N}} \dfrac{n^2 - n + 4}{n+1}\, x^n, 
       \quad \sum_{n \in \mathbb{N}} \dfrac{x^n}{2n-1}. $$

\noindent\Large{\underline{\textbf{Exercise 3}}}
\normalsize{}\\
\noindent Expand the following functions into power series at the origin (Maclaurin series).
    $$ f_1 : x \longmapsto \dfrac{1}{(3+x)^2}, \qquad 
       f_2 : x \longmapsto x\ln\!\left(x + \sqrt{x^2+1}\right), \qquad 
       f_3 : x \longmapsto \arctan\!\left(\dfrac{1-x^2}{1+x^2}\right). $$

\noindent\Large{\underline{\textbf{Exercise 4}}}
\normalsize{}\\
\noindent Use the definition of Taylor series to find the Taylor series, 
centered at $x_0$ for the function $f$ in the following cases :
    $$  f_1 : x \longmapsto \cos x,\ \ x_0 = \dfrac{\pi}{4}; \qquad 
        f_2 : x \longmapsto e^x,\ \ x_0 = 1; \qquad 
        f_3 : x \longmapsto \sqrt{x},\ \ x_0 = 4.   $$

\noindent\Large{\underline{\textbf{Exercise 5}}}
\normalsize{}\\
\begin{enumerate}
    \item Solve in $\mathbb{R}$ the equation $\displaystyle\sum_{n=0}^{+\infty} (3n+1)^2\, x^n = 0$.
    \item Using integration by parts to establish equality
        $$\iint_{[0,1]^2} xy e^{xy}\, dx\, dy = e - 1 - \sum_{n=1}^{+\infty} \dfrac{1}{n\,(n!)}. $$
\end{enumerate}

\noindent\Large{\underline{\textbf{Exercise 6}}}
\normalsize{}\\
\noindent Let $(a_n)_{n \in \mathbb{N}}$ the sequence defined by $a_0 = 1$ and 
$2a_{n+1} = \displaystyle\sum_{k=0}^{n} \mathrm{C}_n^k\, a_k\, a_{n-k}$, 
we put $f : x \longmapsto \displaystyle\sum_{n=0}^{+\infty} \dfrac{a_n}{n!} x^n$.
\begin{enumerate}
    \item Show that $\forall n\in\N,\ a_n \leq n!$. Deduce that the radius of convergence of 
          $\displaystyle\sum_{n \in \mathbb{N}} \dfrac{a_n}{n!} x^n$ is strictly positive.
    \item Calculate $f'$ as a function of $f$, deduce $f$.
    \item Deduce $a_n$ as function of $n$.
\end{enumerate}
